% ---------------------------------------------------------------------------
% Preambolo: pacchetti, stile, macro del progetto.
% Nessun pacchetto esotico: tutto quello che serve sta in una TeX Live standard.
% ---------------------------------------------------------------------------

\usepackage[utf8]{inputenc}
\usepackage[T1]{fontenc}
\IfFileExists{lmodern.sty}{\usepackage{lmodern}}{}
\usepackage[italian]{babel}
\usepackage{microtype}

% Tolleranza tipografica: i nomi di file in tt non si sillabano, e senza un po'
% di elasticita' finirebbero fuori dal margine.
\tolerance=1200
\emergencystretch=3em
\hbadness=10000

\usepackage[a4paper,top=2.6cm,bottom=2.6cm,left=2.7cm,right=2.7cm]{geometry}
\usepackage{amsmath,amssymb}
\usepackage{graphicx}
\usepackage{xcolor}
\usepackage{booktabs}
\usepackage{multirow}
\usepackage{array}
\usepackage{longtable}
\usepackage{enumitem}
\usepackage{float}
\usepackage[font=small,labelfont=bf]{caption}
\usepackage{subcaption}
\usepackage{fancyhdr}
\usepackage{listings}
\usepackage{tcolorbox}
\usepackage{siunitx}
\usepackage{pgfplots}
\pgfplotsset{compat=1.16}
\usepackage{pgfplotstable}
\usepackage[hidelinks]{hyperref}
\usepackage[italian]{cleveref}

% ---------------------------------------------------------------------------
% Unita' di misura
% ---------------------------------------------------------------------------
\sisetup{
  detect-weight       = true,
  detect-family       = true,
  group-digits        = integer,
  group-minimum-digits= 4,
  output-decimal-marker = {,},
  per-mode            = symbol
}
\DeclareSIUnit{\flop}{FLOP}
\DeclareSIUnit{\gflops}{GFLOPS}
\DeclareSIUnit{\gib}{GiB}
\DeclareSIUnit{\mib}{MiB}
\DeclareSIUnit{\kib}{KiB}

% ---------------------------------------------------------------------------
% Intestazioni
% ---------------------------------------------------------------------------
\pagestyle{fancy}
\fancyhf{}
\fancyhead[L]{\small\nouppercase{\leftmark}}
\fancyhead[R]{\small\thepage}
\renewcommand{\headrulewidth}{0.4pt}

% ---------------------------------------------------------------------------
% Notazione ricorrente del progetto
% ---------------------------------------------------------------------------
\newcommand{\Mat}[1]{\ensuremath{#1}}
\newcommand{\mloc}{\ensuremath{m_{\text{loc}}}}
\newcommand{\nloc}{\ensuremath{n_{\text{loc}}}}
\newcommand{\Prow}{\ensuremath{P_r}}
\newcommand{\Pcol}{\ensuremath{P_c}}
\newcommand{\Ai}{\ensuremath{I}}                  % intensita' aritmetica
\newcommand{\code}[1]{\texttt{#1}}
\newcommand{\file}[1]{\texttt{#1}}
\newcommand{\kernelname}[1]{\texttt{#1}}

% ---------------------------------------------------------------------------
% Marcatori dei dati ancora da raccogliere.
%
% Sono volutamente VISIBILI: questa relazione viene compilata mentre le
% campagne sono in corso, e un segnaposto che non si vede e' un segnaposto che
% finisce in consegna. Per la versione finale basta commentare la riga
% \todovisibiletrue qui sotto: i marcatori spariscono dal PDF e restano nel
% sorgente.
% ---------------------------------------------------------------------------
\newif\iftodovisibile
\todovisibiletrue          % <-- commentare questa riga per la consegna finale

\definecolor{todorosso}{HTML}{B3261E}
\definecolor{todogiallo}{HTML}{FFF4E5}
\definecolor{grigiodato}{HTML}{9E9E9E}

% Nota a margine da risolvere
\newcommand{\TODO}[1]{%
  \iftodovisibile
    \textcolor{todorosso}{\textbf{[DA FARE: #1]}}%
  \fi
}

% Singolo numero non ancora misurato, dentro una tabella
\newcommand{\dato}{\textcolor{grigiodato}{--}}

% Blocco di testo che dipende da misure non ancora disponibili
\newenvironment{daverificare}
  {\begin{tcolorbox}[colback=todogiallo,colframe=todorosso,boxrule=0.5pt,
                     left=6pt,right=6pt,top=4pt,bottom=4pt,
                     title=\small\textbf{Da completare con i dati della campagna}]}
  {\end{tcolorbox}}

% Ipotesi interpretativa non ancora confermata da un profilo
\newenvironment{ipotesi}[1]
  {\begin{tcolorbox}[colback=blue!3,colframe=blue!35,boxrule=0.5pt,
                     left=6pt,right=6pt,top=4pt,bottom=4pt,
                     title=\small\textbf{Ipotesi: #1}]}
  {\end{tcolorbox}}

% Segnaposto grafico: riquadro delle dimensioni della figura definitiva
\newcommand{\figuradafare}[2][0.72\linewidth]{%
  \begin{tcolorbox}[colback=todogiallo,colframe=todorosso,boxrule=0.6pt,
                    width=#1,halign=center,valign=center,height=5.2cm]
    \textcolor{todorosso}{\textbf{Figura da generare}}\\[2pt]
    \small #2
  \end{tcolorbox}%
}

% ---------------------------------------------------------------------------
% Listati
% ---------------------------------------------------------------------------
\definecolor{codegrigio}{HTML}{F5F5F5}
\definecolor{codecommento}{HTML}{4C7A34}
\definecolor{codekeyword}{HTML}{1A4F9C}

\lstdefinestyle{scpa}{
  backgroundcolor=\color{codegrigio},
  basicstyle=\ttfamily\footnotesize,
  keywordstyle=\color{codekeyword}\bfseries,
  commentstyle=\color{codecommento}\itshape,
  stringstyle=\color{todorosso},
  numbers=none,
  frame=single,
  framesep=4pt,
  rulecolor=\color{black!25},
  breaklines=true,
  breakatwhitespace=true,
  showstringspaces=false,
  columns=fullflexible,
  keepspaces=true,
  captionpos=b,
  aboveskip=1.0em,
  belowskip=1.0em,
  literate={à}{{\`a}}1 {è}{{\`e}}1 {é}{{\'e}}1 {ì}{{\`i}}1
           {ò}{{\`o}}1 {ù}{{\`u}}1 {°}{{$^\circ$}}1
}
\lstset{style=scpa}
\lstdefinelanguage{cuda}[ANSI]{C++}{
  morekeywords={__global__,__device__,__shared__,__restrict__,__syncthreads,
                __shfl_down_sync,blockIdx,threadIdx,blockDim,cudaMalloc,
                cudaMemcpy,cudaEventRecord,scalar_t}
}
\lstdefinelanguage{shell}{
  morekeywords={make,mpirun,nvcc,module,ompi_info,nsys,ncu,nvidia-smi},
  sensitive=true,
  morecomment=[l]{\#}
}
\renewcommand{\lstlistingname}{Listato}

% ---------------------------------------------------------------------------
% Stile comune dei grafici
% ---------------------------------------------------------------------------
\pgfplotsset{
  scpa/.style={
    width=0.78\linewidth,
    height=6cm,
    grid=both,
    grid style={line width=.1pt, draw=black!12},
    major grid style={line width=.2pt,draw=black!30},
    tick align=outside,
    legend cell align=left,
    legend style={font=\footnotesize,draw=black!25},
    label style={font=\small},
    tick label style={font=\footnotesize},
    every axis plot/.append style={line width=1pt, mark size=2pt}
  }
}
